\item Hace décadas se creía que los grandes dinosaurios herbívoros, como el Apatosaurus y el Brachiosaurus, habitualmente caminaban sobre el fondo de los lagos, y sus largos cuellos les permitían tomar aire de la superficie para respirar. El Brachiosaurus tenía sus orificios nasales en la parte superior de su cabeza. En 1977, Knut Schmidt-Nielsen puntualizó que para tal criatura sería muy difícil el proceso respiratorio. Como un simple modelo, considere una muestra de $10.0 \text{ L}$ de aire a presión absoluta de $2.00 \text{ atm}$, con densidad de $2.40 \text{ kg/m}^3$, localizada en la superficie de un lago de agua dulce. Encuentre el trabajo requerido para transportarla a una profundidad de $10.3 \text{ m}$, permaneciendo constantes su temperatura, volumen y presión. Este requerimiento de energía es mayor que la energía que puede obtenerse mediante el metabolismo de alimentos con el oxígeno en esa cantidad de aire.
